% !TeX root = ../main.tex
% Add the above to each chapter to make compiling the PDF easier in some editors.

\chapter{Conclusion}\label{chapter:conclusion}

This thesis presented a method for improving camera pose and intrinsic estimation from casual, unlabelled video by combining AnyCam's self-supervised pose prediction with AnyCalib's model-agnostic calibration pipeline.

Two contributions were developed. First, the \textbf{AnyCam~$\times$~AnyCalib integration} replaces AnyCam's 32-candidate focal length system with AnyCalib's continuous calibration, injecting accurate intrinsics into the pose head via a learned focal embedding. Second, the \textbf{Multi-Frame Calibration Transformer (MCT)} enforces calibration consistency across video frames by aggregating AnyCalib's backbone features at four resolution scales via cross-frame self-attention before decoding, preserving geometric structure that would be lost in scalar averaging. Both contributions are trained end-to-end using self-supervised flow reprojection losses combined with a calibration anchor, with only 7.5\% of parameters trainable.

Evaluation on MPI Sintel, TUM-RGBD, and KITTI demonstrated consistent improvements over both baselines. Translation direction error---AnyCam's primary weakness---improved by up to 39.5\% (median) on Sintel. Calibration accuracy improved by up to 32.5\% MAPE relative to AnyCalib on TUM-RGBD. The method satisfies all key requirements identified in the related work analysis (Table~\ref{tab:related-work-comparison}): label-free training, feed-forward inference, multi-frame calibration consistency, dynamic scene robustness, and model-agnostic calibration.

The staged training approach also provides a general template for integrating pretrained vision models into self-supervised pipelines: keep large backbones frozen, introduce lightweight components incrementally, and train them jointly once each has been individually validated.
